% ─────────────────────────────────────────────────────────────────────────────
%  §17  Runtime Profiles and Execution Context
%
%  Status: NORMATIVE SPECIFICATION — design agreed 2026-08-17 between
%          Gert Core Team and SQL Live-Site Operations (gert-sqllivesite).
%          Source of truth: .squad/decisions.md (Rev 4 Final Acknowledgment,
%          Barbara arch-eval 2026-08-16, David integration-critique 2026-08-16,
%          Don ground-truth 2026-08-16 and Slice-1 report 2026-08-17,
%          Ken Slice-2 report 2026-08-17).
%
%  This section is a deliverable for the SQL Live-Site Operations team.
%  They MUST NOT annotate tool definitions with AllowedEnvironments values
%  until §\ref{sec:rp-context-vocab} is published.
%
%  Implementation status of every normative claim is marked explicitly:
%    [IMPL]   — implemented and wired in production as of 2026-08-17.
%    [SPEC]   — specified here but not yet implemented or enforced.
%               An external team reading this section MUST treat [SPEC]
%               items as design commitments, not current behavior.
%
%  The key words MUST, MUST NOT, REQUIRED, SHALL, SHALL NOT, SHOULD,
%  SHOULD NOT, RECOMMENDED, MAY, and OPTIONAL are used in the RFC 2119
%  sense throughout this section.
% ─────────────────────────────────────────────────────────────────────────────

\chapter{Runtime Profiles and Execution Context}
\label{ch:runtime-profiles}

A \emph{runtime profile} is a named, versioned YAML document that parameterizes
how a Gert run executes after the tool-definition binding phase is complete.
It declares the execution context (which class of host the run occurs on),
the attendance mode (whether an active approver is present), the approval
scope, transport opt-ins, and per-tool endpoint overrides.

This chapter defines the profile document schema (\S\ref{sec:rp-schema}),
the closed vocabulary of context identifiers (\S\ref{sec:rp-context-vocab}),
attendance declaration (\S\ref{sec:rp-attendance}), profile flatness
(\S\ref{sec:rp-flatness}), the composition of profiles with the
package-map mechanism (\S\ref{sec:rp-composition}), action classification
(\S\ref{sec:rp-classification}), the orthogonality of approval and
classification (\S\ref{sec:rp-approval-classification}), governance
composition rules (\S\ref{sec:rp-governance-composition}),
test-context transport rules (\S\ref{sec:rp-test-transport}),
tiered preflight (\S\ref{sec:rp-preflight}), and the current
implementation status (\S\ref{sec:rp-impl-status}).

% ─────────────────────────────────────────────────────────────────────────────
\section{Profile Document Schema}
\label{sec:rp-schema}
% ─────────────────────────────────────────────────────────────────────────────

A runtime profile document MUST conform to the following shape:

\begin{minted}{yaml}
apiVersion: runtime-profile/v1
id: <profile-id>
context: <one of five values; see §17.2>
attendance: attended | unattended
approval:
  scope:
    allow_read: bool
    allow_mutating: bool
    allow_destructive: bool
  legacy_unspecified_policy: allow | prompt
transport:
  allow_subprocess_in_test: bool
tools:
  <tool-id>:
    endpoint: <url>
\end{minted}

\subsection{Field-by-Field Rules}

\paragraph{\texttt{apiVersion}.} MUST be the literal string
\texttt{runtime-profile/v1}. A document whose \texttt{apiVersion} is absent or
does not match this value MUST be rejected at parse time.

\paragraph{\texttt{id}.} A non-empty string that uniquely identifies the profile
within the project.  [SPEC]

\paragraph{\texttt{context}.} REQUIRED.  One of the five canonical values defined
in \S\ref{sec:rp-context-vocab}.  A profile whose \texttt{context} field is absent
or holds an unrecognized value MUST be rejected at parse time.  [SPEC]

\paragraph{\texttt{attendance}.} REQUIRED.  One of \texttt{attended} or
\texttt{unattended}.  Attendance MUST NOT be inferred from ambient state such as
TTY presence, the value of \texttt{CI} environment variables, or any other runtime
signal.  It MUST be declared explicitly in the profile.  See
\S\ref{sec:rp-attendance}.  [SPEC]

\paragraph{\texttt{approval.scope}.} Controls which action classifications the
profile permits.  All three sub-fields are boolean.  Absent means \texttt{false}.
[SPEC]

\paragraph{\texttt{approval.legacy\_unspecified\_policy}.} Controls the migration
UX for actions whose \texttt{classification} and \texttt{requires-approval}
fields are both unspecified. \texttt{allow} suppresses the migration PROMPT only.
It MUST NOT relax retry, idempotency, or late-result behavior.  See
\S\ref{sec:rp-approval-classification}.  [SPEC]

\paragraph{\texttt{transport.allow\_subprocess\_in\_test}.} Boolean; default
\texttt{false}.  Governs MCP subprocess transport in the \texttt{test} context.
See \S\ref{sec:rp-test-transport}.  [SPEC]

\paragraph{\texttt{tools.<tool-id>.endpoint}.} Per-tool endpoint override.  A
profile MAY override the endpoint URL for a named tool.  A profile MUST NOT
override the transport mode (see \S\ref{sec:rp-composition}).  [SPEC]

% ─────────────────────────────────────────────────────────────────────────────
\section{Execution Context Vocabulary}
\label{sec:rp-context-vocab}
% ─────────────────────────────────────────────────────────────────────────────

\subsection{Canonical Context Identifiers}
\label{subsec:rp-context-values}

The set of valid \texttt{context} values is \textbf{closed}.  There are exactly
five.  Implementations MUST NOT accept any value not in this set.

\begin{center}
\begin{tabular}{ll}
\toprule
\textbf{Value} & \textbf{Meaning} \\
\midrule
\texttt{cli-operator}     & Interactive terminal session driven by a human operator \\
\texttt{vscode-operator}  & VS Code host (human-driven or autonomous agent in VS Code) \\
\texttt{ci}               & Continuous-integration pipeline \\
\texttt{headless-server}  & Autonomous server process (daemon, job, background agent) \\
\texttt{test}             & Local test execution \\
\bottomrule
\end{tabular}
\end{center}

\noindent These values describe the \textbf{execution host} only.  They MUST NOT
encode:
\begin{itemize}[noitemsep]
  \item incident state (e.g., ``this is a live-site incident'');
  \item approval policy (e.g., ``approval required''); or
  \item attendance mode (e.g., ``a human is watching'').
\end{itemize}

Approval policy and attendance are declared as separate top-level fields (see
\S\ref{sec:rp-schema} and \S\ref{sec:rp-attendance}).

\subsection{Correspondence with \texttt{AllowedEnvironments}}
\label{subsec:rp-allowed-environments}

\textbf{Answer to counterparty question (SQL Live-Site Operations, 2026-08-17):}
The value populated in a tool definition's \texttt{AllowedEnvironments} field
MUST correspond to one or more of the canonical context identifiers defined in
\S\ref{subsec:rp-context-values}.  Specifically:

\begin{quote}
  The \texttt{AllowedEnvironments} value on a tool definition corresponds to the
  \texttt{context} field of the selected runtime profile.  A Tier~0 preflight
  check (see \S\ref{sec:rp-preflight}) fires the error code
  \texttt{config/no-binding-for-profile} when the active profile's context is
  not listed in the tool's \texttt{AllowedEnvironments}.
\end{quote}

\noindent This is a normative requirement, not a convention.  Tool definitions MUST
use the identifiers from \S\ref{subsec:rp-context-values} verbatim — no aliases,
no compound values, no free-form strings.

\begin{tcolorbox}[colback=yellow!10, colframe=orange!60, title=Implementation Status]
  \texttt{AllowedEnvironments} is declared in the schema
  (\texttt{pkg/schema/tool.go}).  Runtime enforcement is \textbf{not yet wired}.
  The Tier~0 preflight that evaluates this field against the active profile's
  context is specified here and will land in Phase~1.  [SPEC]
\end{tcolorbox}

% ─────────────────────────────────────────────────────────────────────────────
\section{Attendance Declaration}
\label{sec:rp-attendance}
% ─────────────────────────────────────────────────────────────────────────────

Attendance is the property that governs whether an active human approver is
present during execution.  It is \textbf{orthogonal} to context.

\begin{itemize}[noitemsep]
  \item \texttt{context: vscode-operator} with \texttt{attendance: unattended}
        is a valid and meaningful combination.  It describes an autonomous agent
        process running inside the VS Code host with no human monitoring the
        approval queue.
  \item \texttt{context: cli-operator} with \texttt{attendance: unattended}
        is also valid.  It describes a script running in a terminal that nobody
        is actively watching.
\end{itemize}

\textbf{Operator contexts MUST NOT auto-imply attended.}  A profile whose
\texttt{context} is \texttt{cli-operator} or \texttt{vscode-operator} and whose
\texttt{attendance} field is absent MUST be rejected; the field is required (see
\S\ref{sec:rp-schema}).

\textbf{Edge cases (normative):}

\begin{itemize}[noitemsep]
  \item A profile that declares \texttt{attendance: attended} when no TTY is
        available MUST fail at Tier~0 preflight with error code
        \texttt{config/attendance-mismatch}.  A profile that promises an
        approver who cannot be reached is a configuration error, not a
        runtime transient.  [SPEC]
  \item A profile that declares \texttt{attendance: unattended} with
        \texttt{context: cli-operator} is legal.  It receives unattended
        governance rules.
\end{itemize}

\begin{tcolorbox}[colback=blue!5, colframe=blue!40, title=Non-Normative Note — Current Behavior]
  Current Gert (as of commit \texttt{c810b96}) infers attendance from
  \texttt{TTYOutput}, which is hardcoded to \texttt{true} in
  \texttt{cmd/gert/run.go}.  No \texttt{isatty()} call is made.  As a result,
  any pipeline job that triggers an approval gate will block on stdin
  indefinitely.  This is a known latent bug; profile-declared attendance
  replaces TTY inference and is the fix.  Until profile-driven attendance is
  implemented ([SPEC]), any tool that declares
  \texttt{requires-approval: true} MUST NOT be used in unattended contexts.
\end{tcolorbox}

% ─────────────────────────────────────────────────────────────────────────────
\section{Profile Flatness}
\label{sec:rp-flatness}
% ─────────────────────────────────────────────────────────────────────────────

Runtime profiles are flat.  A profile document MUST NOT contain an
\texttt{extends}, \texttt{inherit}, \texttt{parent}, or any other field that
causes values to be merged from another profile document at load time.

Per-tool overrides within a single profile document are permitted (see
\S\ref{sec:rp-schema}, \texttt{tools.<tool-id>.endpoint}).  These are
overrides within the same document, not cross-document inheritance.

The rationale for this constraint is auditability: the effective configuration
of any run MUST be readable from a single document without recursive
resolution.

% ─────────────────────────────────────────────────────────────────────────────
\section{Profile and Package-Map Composition}
\label{sec:rp-composition}
% ─────────────────────────────────────────────────────────────────────────────

Gert accepts two flags that govern tool binding:

\begin{itemize}[noitemsep]
  \item \texttt{--package-map}: governs \textbf{which tool-definition file
        (YAML)} is loaded for each \texttt{toolRef}.  It operates at the
        catalog/registry layer.  It determines which
        \texttt{*.tool.yaml} on disk backs each logical tool name.
  \item \texttt{--profile}: parameterizes execution \textbf{after} YAML
        selection: context, attendance, endpoint URLs, authentication, and
        per-tool endpoint overrides within the selected tool definitions.
\end{itemize}

These two mechanisms compose; they do not compete.  The composition rules are:

\begin{enumerate}
  \item \texttt{--package-map} wins at the YAML-selection layer.  It
        determines which tool-definition file backs each \texttt{toolRef}.
        The profile's transport configuration is applied to the
        \textbf{already-selected} (post-package-map) tool definition.

  \item The profile-aware binding resolver MUST consume
        \texttt{plan.Tools} — the post-catalog, post-package-map resolved
        tool set — as its input.  It MUST NOT re-resolve \texttt{toolRef}
        declarations independently.  A mismatch between the resolved
        tool-definition's transport mode and the profile's transport
        requirements MUST be surfaced as a Tier~0 error (see
        \S\ref{sec:rp-preflight}); it MUST NOT be silently overridden.
        [SPEC]

  \item A profile MUST NOT silently rewrite a tool definition's transport
        mode.  Different physical implementations (e.g., a mock subprocess
        and a production HTTP endpoint) are separate tool definitions
        selected via \texttt{--package-map}.

  \item If \texttt{--package-map} redirects a \texttt{toolRef} to a mock
        package (transport mode: \texttt{native}) while the active profile
        expects \texttt{mcp-http}, the Tier~0 check reports
        \texttt{config/no-binding-for-profile}.  This is the correct
        behavior: the operator explicitly composed incompatible choices.
\end{enumerate}

\begin{tcolorbox}[colback=blue!5, colframe=blue!40, title=Non-Normative Note — Current State]
  \texttt{--package-map} is implemented and wired.  [IMPL].  The
  profile-aware binding resolver described in rules (2)–(4) is not yet
  implemented.  [SPEC]
\end{tcolorbox}

% ─────────────────────────────────────────────────────────────────────────────
\section{Action Classification}
\label{sec:rp-classification}
% ─────────────────────────────────────────────────────────────────────────────

\subsection{Governing Principle}

\begin{tcolorbox}[colback=red!5, colframe=red!50, title=Governing Principle — Classification]
  Classification is an \textbf{opt-in to reduced friction}.  Absence of a
  classification value MUST NEVER grant additional execution rights, enable
  automatic retry, or relax late-result handling.  Conservative behavior is
  always the default.
\end{tcolorbox}

\noindent This principle applies without exception.  Any implementation that
provides additional permissions or relaxes safety constraints in the absence of a
classification value violates this specification.

\subsection{Field Location and Type}

Classification is a \textbf{per-action} field on \texttt{ToolAction}, not a
per-tool field on \texttt{ToolGovernance}:

\begin{minted}{go}
// pkg/schema/tool.go
type ToolAction struct {
    // ...
    Classification *string `yaml:"classification,omitempty" json:"classification,omitempty"`
}
\end{minted}

A pointer type is used so that absence (\texttt{nil}) is distinguishable from
the explicit value \texttt{"unspecified"}.  [IMPL]

\subsection{Valid Classification Values}

The set of valid classification values is closed.  Implementations MUST reject
any value not in this set at parse time.  The validation check
\texttt{validateActionClassifications()} in \texttt{internal/tool/scan.go}
enforces this.  [IMPL]

\begin{center}
\begin{tabular}{lp{9.5cm}}
\toprule
\textbf{Value} & \textbf{Meaning} \\
\midrule
\texttt{read-only}    & The action has no side effects observable outside the
                         calling context.  Retry is PERMITTED only when the
                         action is also marked explicitly idempotent
                         (see \S\ref{sec:rp-late-result}). \\
\texttt{mutating}     & The action changes external state in a reversible or
                         non-destructive manner. \\
\texttt{destructive}  & The action deletes, overwrites, or permanently alters
                         external state. \\
\texttt{unspecified}  & Explicit declaration that the classification is not
                         known.  This is a permanent classification, not a
                         transitional placeholder.  It receives the same
                         conservative handling as absent. \\
\midrule
\emph{absent (nil)}   & Treated identically to \texttt{unspecified}. \\
\bottomrule
\end{tabular}
\end{center}

\noindent \textbf{\texttt{unspecified} is permanent, not transitional.}  Authors
who cannot determine whether an action is safe to retry MUST use
\texttt{unspecified} or leave the field absent.  Neither value implies that the
author intends to add classification later.

% ─────────────────────────────────────────────────────────────────────────────
\section{Approval and Classification Are Orthogonal}
\label{sec:rp-approval-classification}
% ─────────────────────────────────────────────────────────────────────────────

\begin{tcolorbox}[colback=red!5, colframe=red!50, title=Most Important Section]
  Approval and classification answer different questions and MUST NOT be derived
  from each other.  This section encodes the correction accepted from SQL
  Live-Site Operations in Round~4 of the design negotiation.
\end{tcolorbox}

\paragraph{\texttt{requires-approval} answers:} Is active approval required
before this action is invoked?

\paragraph{\texttt{classification} answers:} What are this action's side effects,
and is automatic retry or late-result handling safe?

\noindent These are independent properties.  A legacy action MAY suppress
approval while being mutating or destructive.  A read-only action MAY require
approval for audit purposes.

\subsection{The \texttt{requires-approval} Tri-State}
\label{subsec:rp-tristate}

\texttt{ToolGovernance.RequiresApproval} is a pointer to bool (\texttt{*bool}).
[IMPL]

\begin{center}
\begin{tabular}{lll}
\toprule
\textbf{Go value} & \textbf{YAML}                     & \textbf{Meaning} \\
\midrule
\texttt{nil}      & absent, or governance block present but field not authored
                  & Unspecified — approval routing defers to context rules \\
\texttt{\&false}  & \texttt{requires-approval: false}
                  & Explicit legacy approval opt-out \\
\texttt{\&true}   & \texttt{requires-approval: true}
                  & Approval required \\
\bottomrule
\end{tabular}
\end{center}

\noindent Prior to commit \texttt{a2e7db0}, \texttt{RequiresApproval} was a
plain \texttt{bool}.  This made \texttt{nil} (unspecified) and \texttt{\&false}
(explicit opt-out) indistinguishable after YAML unmarshal.  The pointer change
corrects this.

\subsection{Non-Coercion Rule}

\texttt{requires-approval: false} (\texttt{\&false}) MUST NEVER assign, imply,
or coerce \texttt{classification: read-only}.

\textbf{Rationale:} The Gert Core Team accepted the following correction from
SQL Live-Site Operations (Round~1, 2026-08-17): an earlier Rev~3 draft stated
that explicit \texttt{requires-approval: false} was ``equivalent to
\texttt{classification: read-only}.''  That statement was wrong.  Coercing
\texttt{\&false} to \texttt{read-only} would have granted retry eligibility and
read-only late-result handling to actions that may be destructive — a safety
regression.

\textbf{Published semantics for the combination
\texttt{RequiresApproval = \&false, Classification = nil}:}

\begin{center}
\begin{tabular}{ll}
\toprule
\textbf{Dimension} & \textbf{Behavior} \\
\midrule
Approval routing    & Legacy opt-out preserved; gate does not fire. \\
Classification      & Remains \texttt{unspecified}. \\
Retry               & MUST NOT retry automatically. \\
Late/lost result    & INDETERMINATE; halt; require verification. \\
\texttt{legacy\_unspecified\_policy: allow} &
    Suppresses migration prompts only.  MUST NOT relax retry,
    idempotency, or late-result behavior. \\
\bottomrule
\end{tabular}
\end{center}

\subsection{Behavior of \texttt{unspecified} by Context}
\label{subsec:rp-unspecified-by-context}

When \texttt{classification} is \texttt{nil} or \texttt{"unspecified"}, the
following rules govern invocation:

\begin{center}
\begin{longtable}{p{4.5cm}p{9cm}}
\toprule
\textbf{Context / situation} & \textbf{Normative behavior} \\
\midrule
Interactive operator (attended) &
    Approval gate fires; active confirmation REQUIRED before invocation. \\
CI / headless (unattended) &
    Invocation is DENIED unless explicitly permitted by
    \texttt{legacy\_unspecified\_policy: allow} in the profile. \\
Test &
    Auto-approve, but only for deterministic transport bindings
    (see \S\ref{sec:rp-test-transport}). \\
Retry after failure &
    MUST NOT retry automatically; regardless of context. \\
Lost or late result &
    Mark INDETERMINATE; halt; require operator verification. \\
\bottomrule
\end{longtable}
\end{center}

% ─────────────────────────────────────────────────────────────────────────────
\section{Late and Lost Result Semantics by Classification}
\label{sec:rp-late-result}
% ─────────────────────────────────────────────────────────────────────────────

When a response is not received within the expected window (transport timeout,
disconnection, or process crash), the following rules apply per classification:

\begin{center}
\begin{longtable}{p{3cm}p{10cm}}
\toprule
\textbf{Classification} & \textbf{Normative behavior} \\
\midrule
\texttt{read-only} &
    Retry is PERMITTED only when the action is explicitly marked idempotent
    (step-level \texttt{retry.idempotent: true}).  If not marked idempotent,
    fail normally — do not retry. \\
\texttt{mutating} &
    Mark INDETERMINATE; halt; require operator verification.
    MUST NOT retry automatically. \\
\texttt{destructive} &
    Mark INDETERMINATE; halt; require operator verification.
    MUST NOT retry automatically. \\
\texttt{unspecified} \emph{or absent} &
    No automatic retry.  If completion cannot be established, halt
    and require operator verification. \\
\bottomrule
\end{longtable}
\end{center}

\noindent \texttt{RequiresApproval} has no bearing on any of these rules.
Retry and late-result behavior are derived solely from \texttt{classification}
and the step-level \texttt{retry} block.  No code path in the current
implementation couples these two properties; the orthogonality is structural
(they inhabit separate types: \texttt{ToolGovernance} and
\texttt{ToolAction} respectively).  [IMPL — structural guarantee; runtime
enforcement of these halt rules is [SPEC].]

% ─────────────────────────────────────────────────────────────────────────────
\section{Runbook-Level vs.\ Per-Action Governance}
\label{sec:rp-governance-composition}
% ─────────────────────────────────────────────────────────────────────────────

A runbook MAY declare a top-level \texttt{governance.require\_approval} field.
The following rules govern how runbook-level and per-action governance compose.

\begin{enumerate}
  \item Runbook-level \texttt{require\_approval} gates the entire run
        uniformly.  It applies to all steps unless overridden at the per-action
        level.

  \item Composition is \textbf{monotone-increasing}.  Runbook-level
        \texttt{require\_approval: false} MUST NOT override or suppress a
        per-action \texttt{requires-approval: true}.

  \item Formally, the effective approval requirement for a step is:

        \[
          \textit{effective} = \textit{runbook\_level} \;\lor\;
          \textit{tool\_level}
        \]

        Because this is a logical OR, the result can only increase.  Once
        either side requires approval, the effective decision is ``required.''

  \item A runbook-level \texttt{require\_approval} MUST NEVER override or
        suppress per-action classification or per-action late-result behavior.
        Classification is an action-level property and is not influenced by
        runbook-level governance.
\end{enumerate}

\paragraph{Implementation.} Monotone composition is implemented at two sites:
\begin{itemize}[noitemsep]
  \item \texttt{pkg/pkgsubst/pkgsubst.go}:
        \verb|eff.RequireApproval = eff.RequireApproval || subRequireApproval|
        on the substitution path.  [IMPL]
  \item \texttt{internal/governance/builder.go:BuildPolicy()}: only ever sets
        \texttt{requireApproval = true}, never resets it to \texttt{false}.
        [IMPL]
  \item \texttt{internal/governance/evaluator.go}: evaluates
        \texttt{e.requireApproval || step.ToolRequiresApproval}.  [IMPL]
\end{itemize}

Test \texttt{TestApproval\_RunbookFalseCannotSuppressToolTrue} demonstrates
that a runbook with \texttt{require\_approval: false} does not suppress a tool
with \texttt{requires-approval: true}.  [IMPL]

\begin{tcolorbox}[colback=blue!5, colframe=blue!40,
  title=Non-Normative Note — Profile-Driven Approval Gate]
  The current approval gate is selected from TTY state, not from the declared
  profile.  A profile-driven approval gate (one that consults
  \texttt{attendance} and \texttt{approval.scope} from the profile document)
  is specified here but not yet implemented.  [SPEC]
\end{tcolorbox}

% ─────────────────────────────────────────────────────────────────────────────
\section{Test-Context Transport Rules}
\label{sec:rp-test-transport}
% ─────────────────────────────────────────────────────────────────────────────

When the active profile's \texttt{context} is \texttt{test}, the following
transport rules apply at Tier~0 preflight (see \S\ref{sec:rp-preflight}).

\begin{center}
\begin{longtable}{p{3cm}p{10cm}}
\toprule
\textbf{Transport mode} & \textbf{Rule in \texttt{test} context} \\
\midrule
\texttt{native} &
    Always allowed.  No profile opt-in required. \\
\texttt{mcp} (subprocess) &
    Blocked by default.  Allowed only with
    \texttt{transport.allow\_subprocess\_in\_test: true} explicitly set
    in the profile. \\
\texttt{mcp-http} &
    Never allowed.  No override exists.  This prohibition is
    technically enforced at Tier~0 (static check); it cannot be
    bypassed by any profile field. \\
\bottomrule
\end{longtable}
\end{center}

\subsection{Subprocess Opt-In: Governance Assertion, Not Technical Sandbox}
\label{subsec:rp-subprocess-assertion}

Setting \texttt{transport.allow\_subprocess\_in\_test: true} is an
\textbf{auditable governance assertion} by the profile author, not a technical
sandbox guarantee.

The profile author asserts, by setting this field, that the named subprocess
command is a hermetic test double and will not contact production systems or
produce irreversible side effects.

\textbf{Gert does not enforce isolation.}  Specifically:
\begin{itemize}[noitemsep]
  \item No Linux namespace isolation (\texttt{SysProcAttr.Cloneflags} is not
        set).
  \item No seccomp filter.
  \item No network restriction.
  \item No working-directory jail.
  \item No environment scrubbing: the subprocess inherits the full parent
        process environment unconditionally
        (\texttt{mergeEnv()} prepends \texttt{os.Environ()} before
        appending tool-specific variables).
\end{itemize}

The subprocess sees everything the Gert process sees.  This behavior is
confirmed in \texttt{internal/tool/process.go:StartProcess()}.  [IMPL]

By contrast, the \texttt{mcp-http} prohibition in \texttt{test} context
\textbf{is} technically enforced at Tier~0: the static check fires before any
process is started or any network call is made.

The rationale for the subprocess escape hatch: a hermetic local MCP server
(e.g., a fake binary producing deterministic responses) is a legitimate pattern
for integration-testing the MCP transport layer.  Banning subprocess entirely
in \texttt{test} context would foreclose that pattern.  The opt-in is explicit,
named, and appears in the profile, making its use auditable.  HTTP endpoints
have no such legitimate test use case at the Gert layer and are banned
unconditionally.

\paragraph{Error codes.}

\begin{center}
\begin{tabular}{lp{9cm}}
\toprule
\textbf{Code} & \textbf{Message conditions} \\
\midrule
\texttt{config/test-context-non-native-binding}
  & Fires for any non-native, non-opted-in transport in a test-context
    profile.  The \texttt{mcp-http} variant cannot be overridden; the
    \texttt{mcp} variant is suppressible with
    \texttt{allow\_subprocess\_in\_test: true}. \\
\bottomrule
\end{tabular}
\end{center}

% ─────────────────────────────────────────────────────────────────────────────
\section{Tiered Preflight}
\label{sec:rp-preflight}
% ─────────────────────────────────────────────────────────────────────────────

Preflight validation is divided into three tiers.  Higher tiers require
additional I/O; lower tiers are always required.

\subsection{Tier~0 — Static / Offline (Mandatory)}
\label{subsec:rp-preflight-t0}

Tier~0 performs pure config-graph traversal with \textbf{zero I/O}: no
network calls, no subprocess execution, no credential operations, no
filesystem reads beyond the already-loaded plan.  Tier~0 MUST execute before
every run.  It is not optional and cannot be skipped.

Tier~0 MUST be able to answer the question ``can this runbook run in this
context?'' with no credentials and no network access.

The following checks constitute Tier~0:

\begin{enumerate}
  \item Every \texttt{toolRef} resolves to a \texttt{tool.yaml} in a loaded
        package (existing PLAN-010 gate).  [IMPL]
  \item The resolved tool's \texttt{transport.mode} has a binding entry in
        the active runtime profile.  [SPEC]
  \item That binding entry is structurally complete for its transport type:
        \texttt{mcp-http} requires \texttt{endpoint} and
        \texttt{auth.provider}; \texttt{mcp} (subprocess) requires
        \texttt{command}; \texttt{native} requires nothing.  [SPEC]
  \item The referenced auth provider is a known type with required
        configuration fields present.  [SPEC]
  \item Environment variable interpolations (e.g.,
        \texttt{\$\{GERT\_MCP\_ENDPOINT\}}) resolve to non-empty values.
        [SPEC]
  \item \texttt{AllowedEnvironments}, if declared on the tool, includes the
        active profile's \texttt{context} value
        (see \S\ref{subsec:rp-allowed-environments}).  [SPEC]
  \item Test-context transport rules from \S\ref{sec:rp-test-transport}
        pass.  [SPEC]
  \item If the profile declares \texttt{attendance: attended} and no TTY is
        available, emit \texttt{config/attendance-mismatch}.  [SPEC]
\end{enumerate}

The Tier~0 input is \texttt{plan.Tools} — the post-catalog,
post-package-map resolved tool set — not raw \texttt{toolRef} declarations.

\subsection{Tier~1 — Local / No Network (Opt-in)}
\label{subsec:rp-preflight-t1}

Activated with \texttt{--preflight=local}.  Tier~1 runs Tier~0 first, then
additionally checks:

\begin{itemize}[noitemsep]
  \item Subprocess \texttt{command} is on \texttt{PATH} and executable.
  \item \texttt{az} (Azure CLI) is on \texttt{PATH} when an
        \texttt{azure-cli} auth provider is configured.
  \item IMDS (\texttt{169.254.169.254}) responds within 500~ms
        (link-local reachability probe; no token is requested).
  \item Native-transport runbook files exist on disk.
\end{itemize}

[SPEC]

\subsection{Tier~2 — Live / Network (Opt-in)}
\label{subsec:rp-preflight-t2}

Activated with \texttt{--preflight=live}.  Tier~2 runs Tier~0 and Tier~1
first, then additionally performs:

\begin{itemize}[noitemsep]
  \item Token acquisition (actual credential operation against the auth
        provider).
  \item Endpoint probe (HTTP HEAD or equivalent against the declared
        endpoint URL).
\end{itemize}

\texttt{--preflight=live} adds Tier~2 on top of Tier~0 — never instead of it.
Tier~2 failures surface the dynamic error codes \texttt{auth/credential-failure}
and \texttt{transport/endpoint-unreachable}.

[SPEC]

\subsection{Error Taxonomy}
\label{subsec:rp-error-taxonomy}

The static vs.\ dynamic distinction MUST be explicit in operator-facing error
message text.  Operator remediation is completely different for static
configuration errors vs.\ runtime network errors.

\begin{center}
\begin{longtable}{p{3.5cm}cp{7.5cm}}
\toprule
\textbf{Code} & \textbf{Tier} & \textbf{Meaning} \\
\midrule
\texttt{config/tool-unresolved} & 0 &
    \texttt{toolRef} names a tool in no loaded package. \\
\texttt{config/no-binding-for-profile} & 0 &
    Tool has no binding in the active profile.  The error message MUST
    include which profiles DO provide a valid binding, if computable. \\
\texttt{config/binding-incomplete} & 0 &
    Binding exists but required configuration field is absent.  Message
    MUST name the missing field and the unset environment variable, if
    applicable. \\
\texttt{config/attendance-mismatch} & 0 &
    Profile declares \texttt{attended} but no TTY is available. \\
\texttt{config/test-context-non-native-binding} & 0 &
    Test-context profile references a non-native transport (see
    \S\ref{sec:rp-test-transport}). \\
\texttt{auth/credential-failure} & 2 &
    Token acquisition failed.  Message MUST note this is a runtime
    error; the configuration binding is valid. \\
\texttt{transport/endpoint-unreachable} & 2 &
    Endpoint probe failed.  Message MUST note this is a runtime/network
    error; the binding and credentials are valid. \\
\bottomrule
\end{longtable}
\end{center}

[SPEC — all error codes listed; Tier~0 code \texttt{config/tool-unresolved}
corresponds to existing PLAN-010 gate [IMPL]; all others are [SPEC].]

\subsection{\texttt{gert plan} Command}
\label{subsec:rp-gert-plan}

\texttt{gert plan --profile <id> <runbook>} is a non-executing binding table
report that runs Tier~0 preflight and displays per-tool binding status.
Example output:

\begin{minted}{text}
$ gert plan --profile headless-server runbooks/icm-tsg-router.runbook.yaml
  toolRef            transport   auth              endpoint               status
  icm                mcp-http    managed-identity  ${ICM_MCP_ENDPOINT}    bound
  tsg-recommendation mcp-http    managed-identity  ${TSG_MCP_ENDPOINT}    bound
\end{minted}

Exit codes: \texttt{0} = Tier~0 clean; \texttt{1} = Tier~0 failures;
\texttt{2} = Tier~2 failures (only with \texttt{--preflight=live}).
\texttt{--output=json} for CI consumption.

\texttt{gert plan --show-profiles <runbook>} lists all profiles in which every
\texttt{toolRef} has a complete binding.

\textbf{Explicit scope note:} \texttt{gert plan} reports what the tool YAML
declares.  It cannot validate endpoint reachability or token acquisition (those
are Tier~2 runtime probes).  It cannot show profile-overridden transport
configurations until the profile binding resolver lands.

[SPEC]

% ─────────────────────────────────────────────────────────────────────────────
\section{Implementation Status}
\label{sec:rp-impl-status}
% ─────────────────────────────────────────────────────────────────────────────

The table below summarizes the implementation status of every normative
requirement in this chapter as of 2026-08-17.  An external team reading this
specification MUST treat [SPEC] entries as design commitments for Phase~1
implementation, not current behavior.

\begin{center}
\begin{longtable}{p{7cm}cp{4cm}}
\toprule
\textbf{Requirement} & \textbf{Status} & \textbf{Evidence / commit} \\
\midrule
Tri-state \texttt{RequiresApproval} (\texttt{*bool} on
\texttt{ToolGovernance}) &
[IMPL] & \texttt{a2e7db0}; \texttt{tool\_governance\_test.go} \\
Per-action \texttt{classification} field with enum validation &
[IMPL] & \texttt{a2e7db0}; \texttt{scan.go} \\
\texttt{AllowedModes} field on \texttt{ToolGovernance} &
[IMPL — field only] & \texttt{a2e7db0} \\
\texttt{GovernanceEvaluator} wired in production &
[IMPL] & \texttt{c810b96}; \texttt{wire.go:161}, \texttt{run.go} \\
Approval gate fires pre-dispatch (all transports) &
[IMPL] & \texttt{c810b96}; \texttt{engine.go} chokepoint \\
Monotone governance composition (runbook + per-tool OR) &
[IMPL] & \texttt{c810b96}; \texttt{evaluator.go}, \texttt{builder.go} \\
Approval does not affect classification or retry &
[IMPL — structural] & Type separation; no coupling in evaluator \\
\midrule
Profile document schema (\texttt{runtime-profile/v1}) & [SPEC] & Phase~1 \\
Profile-driven approval gate (attendance-aware) & [SPEC] & Phase~1 \\
Declared attendance (replacing TTY inference) & [SPEC] & Phase~1 \\
\texttt{AllowedEnvironments} enforcement at Tier~0 & [SPEC] & Phase~1 \\
\texttt{AllowedModes} enforcement at runtime & [SPEC] & Phase~1 \\
Tier~0 preflight checks (2–8 in \S\ref{subsec:rp-preflight-t0}) & [SPEC] & Phase~1 \\
Tier~1 and Tier~2 preflight & [SPEC] & Phase~1 \\
\texttt{gert plan} command & [SPEC] & Phase~1 \\
Profile binding resolver on post-catalog \texttt{plan.Tools} & [SPEC] & Phase~1 \\
Late-result halt rules by classification & [SPEC] & Phase~1 \\
Profile-driven \texttt{approval.scope} enforcement & [SPEC] & Phase~1 \\
\bottomrule
\end{longtable}
\end{center}
